
%\needspace{4\baselineskip}
\color{uniblue}{\subsection{Дистанционная защита (ДЗ)}}
\color{black}

\begin{enumerate}[label=\arabic{section}.\arabic{subsection}.\arabic{enumi}, labelsep=4pt, leftmargin=0pt, itemindent=57pt, itemsep=0pt, parsep=5pt] % если что можно крутить itemsep
%%%%%%%%%%%%%%%%%%%%%%%%%%%%%%%%%%%%%%%%%%%%%%%%%%%%%%%%%%%%%%%%%%%%%%  ОБЩИЕ СВЕДЕНИЯ %%%%%%%%%%%%%%%%%%%%%%%%%%%%%%%%%%%%%%%%%%%%%%%%%%
\item Общие сведения

\begin{enumerate}[label=\arabic{section}.\arabic{subsection}.\arabic{enumi}.\arabic*, labelsep=4pt, leftmargin=0em, itemindent=65pt, parsep=0pt]

\item
Дистанционная защита используется для защиты энергообъектов от междуфазных КЗ. Защита срабатывает при снижении замера  сопротивления  сети,  т.е.  по  принципу  действия  является  минимальной.  ДЗ  является  защитой с относительной  селективностью  и  выполняет  функции  основной  или  резервной  защиты. Принцип работы функции основан на измерении полного сопротивления. Ступени ДЗ выполнены с независимыми выдержками времени и предназначены для селективного действия при междуфазных коротких замыканиях.

\item 
Основными преимуществами ДЗ являются:
\begin{itemize}
    \item надежная защита определенной части защищаемого объекта от междуфазных КЗ вне зависимости от режимов питающей системы; 
    \item может  применяться  для  обеспечения  дальнего  резервирования.
\end{itemize}

\item 
Селективность защиты обеспечивается введением ступенчатых выдержек времени:
\begin{itemize} 
\item все замыкания в пределах первой зоны (зона охвата первой ступени) отключаются с минимальным временем; 
\item замыкания в пределах 2-3 зоны -- с большим временем.
\end{itemize}

\item
Основные характеристики функционального блока <<ДЗ>>:
\begin{itemize}
    \item минимальный ток пуска реле сопротивления составляет 5\% от $I_{\textsl{ном}}$;
    \item погрешность реле сопротивления не более 5\%;
    \item погрешность по углу при номинальных значениях тока и напряжения не более 5$^\circ$; 
    \item время срабатывания реле сопротивления при угле, равном углу линии, трехкратном токе точной работы и скачкообразном уменьшении напряжения в два раза по отношению к номинальному не более 25 мс;
    \item время возврата реле сопротивления при угле, равном углу линии, трехкратном токе точной работы и скачкообразном увеличении напряжения от 0,1 до 1,0 номинального не более 25 мс;
    \item время срабатывания реле сопротивления при скачкообразном изменении сопротивления на входе от $1,2(X_{\textsl{уст}}/\cos\phi_{\textsl{нагр}})$ до $0,6(X_{\textsl{уст}}/\cos\phi_{\textsl{нагр}})$ при угле максимальной чувствительности не более 25 мс;
    \item \sloppy время возврата реле сопротивления при скачкообразном изменении сопротивления на входе от $0,1(X_{\textsl{уст}}/\cos\phi_{\textsl{нагр}})$ до $1,2(X_{\textsl{уст}}/\cos\phi_{\textsl{нагр}})$ при угле максимальной чувствительности не более 50 мс.
\end{itemize}

\item
Функциональный блок <<ДЗ>>, реализованный в устройстве, включает в себя три ИО полного сопротивления для междуфазных КЗ (канал <<фаза-фаза>>). В его составе реализованы три независимых ступени дистанционной защиты:
\begin{itemize}
    \item <<ДЗ-фф 1 ст.>> (1 ступень ДЗ, измерение сопротивления по контуру <<фаза-фаза>>);
    \item <<ДЗ-фф 2 ст.>> (2 ступень ДЗ, измерение сопротивления по контуру <<фаза-фаза>>);
    \item <<ДЗ-фф 3 ст.>> (3 ступень ДЗ, измерение сопротивления по контуру <<фаза-фаза>>).         
\end{itemize}

\item
Функциональный блок <<ДЗ>> может быть полностью выведен из работы оперативно путем активации сигнала <<ОВ ЗДЗ>>. 

\item 
Устройство контролирует значения следующих величин:
\begin{itemize} 
\item междуфазные  напряжения  ($U_{ab}$,  $U_{bc}$,  $U_{ca}$).  Данные  величины  рассчитываются  как  векторная разность соответствующих фазных напряжений; 
\item междуфазные  токи  ($I_{ab}$,  $I_{bc}$,  $I_{ca}$).  Данные  величины  рассчитываются  как  векторная  разность соответствующих фазных токов. 
\end{itemize} 

\item
Условием  срабатывания  реле  сопротивления  является  снижение  активной  и  реактивной  части измеренного сопротивления ниже соответствующих уставок. Алгоритм расчета определяет три различных сопротивления $Z$ для трех контуров <<фаза-фаза>> каждой ступени ($\dot Z_{ab}$, $\dot Z_{bc}$, $\dot Z_{ca}$). Дистанционная  защита работает таким образом, что измерительные органы ступеней защиты одновременно контролируют все контуры <<фаза-фаза>>. 

\item
Ступени ДЗ <<ДЗ-фф 1 ст.>>, <<ДЗ-фф 2 ст.>> и <<ДЗ-фф 3 ст.>> выполняют  функцию  защиты  от междуфазных  КЗ,  осуществляя контроль петли <<фаза-фаза>>. Ступень ДЗ <<ДЗ-фф 3 ст.>>  предназначена  для выполнения функции дальнего резервирования. 

\item Для правильной работы органа направления к ИО должны подводиться напряжения, рассчитанные по формулам:
% ФОРМУЛЫ №1
\begin{subequations}\label{eq:group1}
    \begin{align}
        \dot U_{ab\textsl{ расч.поляр.}} = \dot U_{1M}(t)\angle{30^\circ}, \\
        \dot U_{bc\textsl{ расч.поляр.}} = \dot U_{1M}(t)\angle{-90^\circ}, \\
        \dot U_{ca\textsl{ расч.поляр.}} = \dot U_{1M}(t)\angle{150^\circ}, 
    \end{align}
\end{subequations}

\begin{eqexpl}[25mm]
    \item{$U_{1}$}вектор прямой последовательности, а
\end{eqexpl}
\begin{equation*}
    \dot U_{1M}(t) = \dot U_{1}(t)+0,125 \cdot \dot U_{\textsl{1памяти}} ,
\end{equation*}

и фазы векторов сопротивления, рассчитанные по формулам:

\begin{subequations}\label{eq:group2}
    \begin{align}
        ang \left( \dot Z_{ab}(t) = \frac{\dot U_{1M}(t)\angle{30^\circ} }{\dot I_a(t) - \dot I_b(t)} \right), \\
        ang \left( \dot Z_{bc}(t) = \frac{\dot U_{1M}(t)\angle{-90^\circ} }{\dot I_b(t) - \dot I_c(t)} \right), \\
        ang \left( \dot Z_{ca}(t) = \frac{\dot U_{1M}(t)\angle{120^\circ} }{\dot I_c(t) - \dot I_a(t)} \right),
    \end{align}
\end{subequations}

\item
Для правильной работы измерительных органов сопротивления дистанционной защиты от аналогового модуля должны приходить реальная и мнимая части полных комплексных сопротивлений $\dot Z_{ab}$, $\dot Z_{bc}$, $\dot Z_{ca}$, рассчитанных по формулам:

\vspace{-\abovedisplayskip}
\begin{subequations}\label{eq:group6}
    \begin{align}
        \dot Z_{ab}(t) = \frac{\dot U_{a}(t) - \dot U_{b}(t)}{\dot I_a(t) - \dot I_b(t)},\\
        \dot Z_{bc}(t) = \frac{\dot U_{b}(t) - \dot U_{c}(t)}{\dot I_b(t) - \dot I_c(t)},\\ 
        \dot Z_{ca}(t) = \frac{\dot U_{c}(t) - \dot U_{a}(t)}{\dot I_c(t) - \dot I_a(t)},       
    \end{align}
\end{subequations}

\item
При выведенной направленности ступени ДЗ имеют многоугольную характеристику срабатывания с охватом нуля (рисунок \ref{dz:img0}). Принцип работы дистанционной защиты с данной  многоугольной характеристикой основан на проверке попадания значений вычисленных сопротивлений в зону срабатывания, которая определяется значениями уставок по сопротивлению (Rср, Xср, Фхс, Фкн). Коэффициент возврата органа построения характеристики равен 1,05. 

%\vspace{2mm}
\begin{figure}[H]
\centering

\includegraphics[width=0.6\textwidth,height=0.6\textheight,keepaspectratio]{\fbpath/301. ДЗ ВВ/_latex/img/1.pdf}
\captionof{figure}{Характеристика срабатывания реле сопротивления ненаправленных ступеней ДЗ}\label{dz:img0}
\end{figure}

\item
Характеристика строится по следующим уставкам:
\begin{itemize} 
    \item уставкой  <<Фхс>>  задается  угол  наклона  характеристики.  В  общем  случае,  угол  наклона характеристики может быть равен углу линии Фл; 
    \item уставкой  <<Rср>>  задаются  активные  сопротивления  (пересечение  правой  и  левой  границ характеристики с осью $R$); 
    \item уставкой  <<Xср>>  задаются  реактивные  сопротивления  (пересечение  верхней  и  нижней  границ характеристики с осью $jX$); 
    \item уставкой  <<Фкн>>  задается  угол,  обеспечивающий  отстройку  от  искажения  замера  на  линиях  с двухсторонним питанием при КЗ через переходное сопротивление; 
    \item для отстройки от срабатывания в нагрузочном режиме, уставками <<Rнг>> и <<Фнг>> задается вырез нагрузки. Учет нагрузочного режима в характеристике срабатывания должен вводиться в работу уставкой <<Вырез\_нагрузки>>.
\end{itemize}

\item
Верхняя сторона характеристики должна четко фиксировать концы защищаемых зон ступеней ДЗ (например, 1 и 2) и быть отстроенной от реактивных (в основном) небольших нагрузок, обуславливающих  значительные реактивные слагающие для 3 ступени. Правая боковая сторона обеспечивает отстройку от рабочих  режимов.  Левая  боковая  сторона  помогает  отстройке  от  мощностей  нагрузок,  передаваемых  к месту включения защиты. Нижняя сторона для органа обеспечивает его работу при близких повреждениях через переходное сопротивление.
\item
Величина угла Фхс (от $\angle{30^\circ}$ до $\angle{89,9^\circ}$) задается в ПО <<ЮНИТ-Сервис>>. В общем случае, угол наклона характеристики может быть равен углу линии Фл. 
\item
В каждой ступени ДЗ предусмотрена возможность выбора направленного режима работы. Работа в направленном режиме осуществляется  с  помощью  органа направления (ОН). Подробное описание органа направления (ОН) представлено в разделе \ref{sec:on} данного руководства.

\item
Для  надежного  срабатывания  защиты  без  учета  направления,  характеристика  смещена  вниз  на величину $0,2 \cdot X$ от оси абсцисс, а уставки по $R$, $X$ и угол наклона должны совпадать с аналогичными для реле сопротивления направленной второй ступени.
\item
В функции дистанционной защиты существует возможность использования подхвата срабатывания первой ступени ДЗ от второй ступени ДЗ. Подхват ступени вводится в работу уставкой SGF9.

\item
Ступени  ДЗ  предусматривают  работу  в  направленном  режиме.  Если  предусмотрен  контроль направленности  при  АУ  (уставка  SGF4  введена  в  работу),  то  наличие  сигнала  <<АУ:  Ввод  ускорения>> позволяет выполнить пуск ступени ДЗ на момент включения выключателя без контроля направления. В каждой ступени ДЗ предусмотрена  возможность  выбора  направленного  режима  работы  контуров. Контроль  направленности  осуществляется  с  помощью  уставки  SGF3.  Для  обеспечения  строгой направленности  действия  ступени,  используется  орган  направления  (ОН).  Направление  действия меняется  в  зависимости  от  положения  уставки  <<Ненаправленная>>,  <<Прямонаправленная>>, <<Обратнонаправленная>>. Характеристики срабатывания прямонаправленных и обратнонаправленных ступеней ДЗ приведены на рис. \ref{dz:img00}.

%\vspace{2mm}
\begin{figure}[H]
\centering
\includegraphics[width=1\textwidth,height=1\textheight,keepaspectratio]{\fbpath/301. ДЗ ВВ/_latex/img/2.pdf}
\captionof{figure}{Характеристики срабатывания реле сопротивления прямонаправленных и обратнонаправленных ступеней ДЗ}\label{dz:img00}
\end{figure}

\item
В  случае  нарушения  устойчивости  параллельной  работы  электростанций  по  линиям электропередачи  возникают  качания  мощности,  сопровождающиеся  значительными  колебаниями  тока, напряжения  и  угла  между  ними.  Т.к.  реле  сопротивления  (PC)  могут  сработать  при  этих  колебаниях, ступени ДЗ дополняются данной блокировкой на время качаний. Чтобы избежать ложного срабатывания ступеней ДЗ в случае возникновения режима качаний, предусмотрен пуск от БК. C помощью уставки SGF7 контролируется  режим  пуска  от  быстродействующих  защит  <<БК:  Срабатывание  БК-б>>  обеспечивающих кратковременный ввод или медленнодействующих защит <<БК: Срабатывание БК-м>>. Во время появления активного уровня <<БК: Срабатывание БК-б>> разрешается срабатывание быстродействующих ступеней ДЗ, с последующей их блокировкой на время присутствия активного уровня сигнала <<БК: Срабатывание БК-м>>. Активный сигнал <<БК: Срабатывание БК-м>> разрешает работу медленнодействующих ступеней и по истечению времени ввода их в работу схема возвращается в исходное состояние.
\item
Уставкой SGF5 в функции предусмотрен выбор режима, от каких ПО будут формироваться пусковые условия работы ступеней ДЗ: <<Без контроля>>, <<С контролем от БК>>, <<С контролем от ОПО>>. При работе с ДЗ существует возможность выбрать способ пуска ступеней защиты:
\begin{itemize}  
    \item с пуском по току (контроль по току осуществляется постоянно); 
    \item с пуском по току и напряжению; 
    \item с пуском по току или по току и напряжению. 
\end{itemize}
\item
С помощью уставки SGF6 контролируется от каких ПО будет осуществлен пуск ступени ДЗ. Работа ДЗ  с  пуском  по  току  и  напряжению  будет  осуществлена  в  том  случае,  когда  хотя  бы  одно  линейное напряжение  становится  меньше  уставки  и  одновременно  ток  в  любой  из  фаз  превышал  порог срабатывания. При низкой чувствительности защиты по уровню напряжения она способна активироваться при превышении током заданной величины ($I_{\textsl{ср}}$). 
\item
При  неисправностях  в  цепях  напряжения  ТН,  напряжение,  подводимое  к  реле  сопротивления, исчезает или резко понижается. В результате этого реле сопротивления, включенные на это напряжение, фиксируют снижение сопротивления и приходят в действие, что приводит к неправильному срабатыванию ДЗ. При неисправном  состоянии  цепей  напряжения  дистанционная  защита  должна  немедленно выводиться из работы. Для исключения ложного срабатывания ступеней ДЗ при неисправностях в цепях ТН, используется блокировка при неисправности цепей напряжения (БНН). При наличии активного сигнала <<БННш: Срабатывание>> и введенной в работу уставке SGF8, ступени ДЗ автоматически блокируются до устранения неисправности в цепях напряжения.
\item
Пусковые  условия  функции  сформируются  только  в  том  случае,  если  отсутствуют  какие-либо блокировки. При  активации  сигнала  пуска,  реле  сопротивления  при  котором  остается  в  сработанном состоянии, запускается таймер Т1 выдержки времени на срабатывание ступени ДЗ. По истечении времени таймера Т1, и отсутствии сигнала блокировки формируется срабатывание ступени ДЗ.

\item
Каждая  ступень  ДЗ  имеет  возможность  оперативного  ускорения.  Оперативное  ускорение  может вводиться:
\begin{itemize} 
    \item при выводе основной защиты оборудования; 
    \item при опробовании оборудования рабочим напряжением; 
    \item при операциях с разъединителями на противоположном конце ВЛ.
\end{itemize}
\item
Для каждой ступени ДЗ предусмотрена возможность ввода автоматического  ускорения при включении выключателя. Функция контроля отсутствия напряжения позволяет ускорять защиту ДЗ при опробовании высоковольтного оборудования. Контроль отсутствия напряжения осуществляется защитой минимального напряжения, цепи которой подключены к ТН, установленному на линии.

\end{enumerate}

%%%%%%%%%%%%%%%%%%%%%%%%%%%%%%%%%%%%%%%%%%%%%%%%%%%%%%%%%%%%%%%%%%%%%%  СТУПЕНИ %%%%%%%%%%%%%%%%%%%%%%%%%%%%%%%%%%%%%%%%%%%%%%%%%%
\item\label{sec:sec2} Ступени дистанционной защиты
\begin{enumerate}[label=\arabic{section}.\arabic{subsection}.\arabic{enumi}.\arabic*, labelsep=4pt, leftmargin=0em, itemindent=65pt, parsep=0pt]

\item
Уставки  <<Вырез  нагрузки>>,  <<Rнг>>,  <<Фнг>>  являются  общими  для  всех  ступеней ДЗ. Общие параметры для настройки функции <<ДЗ>> приведены в таблице \ref{dz:tbl0}.

%================================================== Таблица =====================================
\small
\begin{longtable}{|>{\centering\arraybackslash}m{5.3cm}|>{\centering\arraybackslash}m{3.3cm}|>{\centering\arraybackslash}m{4.2cm}|>{\centering\arraybackslash}m{1.8cm}|>{\centering\arraybackslash}m{1cm}|}
\caption{Общие параметры для настройки функции <<ДЗ>>\hfill\vspace{-0.5\baselineskip}}\label{dz:tbl0}\\ 
\hline
\rowcolor{gray!30}
Параметр (Параметр на ИЧМ) & Условное обозначение на схеме & Значение/ Диапазон & Единица измерения & Шаг \\ 
\hline
\endfirsthead
\caption*{\hspace{3pt}\emph{Продолжение таблицы \ref{dz:tbl0}\hfill\vspace{-0.5\baselineskip}}} \\ % сделано по ГОСТ 2.105 п.6.8.7
\hline
\rowcolor{gray!30}
Параметр (Параметр на ИЧМ) & Условное обозначение на схеме & Значение/ Диапазон & Единица измерения & Шаг \\ 
%\hline
\endhead
%\hline
\endfoot
%\hline
\endlastfoot
%==+t1*LLN0|VV_LVDIS
\centering Учет нагрузочного режима в характеристике срабатывания (Вырез\_нагрузки) & \centering SGF2 & \centering 0 = Не предусмотрено\\1 = Предусмотрено & \centering -- & \centering \arraybackslash -- \\
\hline
\centering Активное сопротивление для отстройки от режима максимальной нагрузки (Rнг) & \centering -- & \centering 0,01 ... 500,00 & \centering Ом & \centering \arraybackslash 0,01 \\
\hline
\centering Угол наклона для отстройки от режима максимальной нагрузки (Фнг) & \centering -- & \centering 5,0 ... 60,0 & \centering град. & \centering \arraybackslash 0,1 \\
\hline
%===t1
\end{longtable}
\normalsize
%================================================== Конец таблицы =====================================

\item
Ввод каждой ступени в работу на этапе параметрирования устройства осуществляется программным переключателем SGF1 <<Ввод\_функции>>. Информация о введенной ступени отображается выходным сигналом <<ДЗ-фф (n)ст.: Ввод>>.
\item
Каждая ступень <<ДЗ>> может быть выведена из работы оперативно путем активации сигнала <<ОВ ДЗ-фф (n)ст.>>, что характеризуется активным состоянием сигнала <<ДЗ-фф (n)ст.: Оперативный вывод>> на выходе функции.

\item
При активном состоянии входного сигнала <<ДЗ-фф(n)ст. на сигн>> действие ступени переводится на сигнал.

\item
Алгоритм работы первой ступени ДЗ по контуру <<фаза-фаза>> (<<ДЗ-фф 1ст.>>) выполнен в соответствии с рисунком \ref{dz:dz1ffimg1}. В таблице \ref{dz:dz1fftbl1} приведены параметры, необходимые для настройки функции <<ДЗ-фф 1ст.>>.

\vspace{3mm}
\begin{figure}[H]
\centering
\includegraphics[width=0.9\textwidth,height=0.9\textheight,keepaspectratio]{\fbpath/301. ДЗ ВВ/_latex/img/dz1ff.pdf}
\captionof{figure}{Функциональная схема <<ДЗ-фф 1ст.>>}\label{dz:dz1ffimg1}
\end{figure}

\small
\begin{longtable}{|>{\centering\arraybackslash}m{5.3cm}|>{\centering\arraybackslash}m{3.3cm}|>{\centering\arraybackslash}m{4.2cm}|>{\centering\arraybackslash}m{1.8cm}|>{\centering\arraybackslash}m{1cm}|}
\caption{Параметры для настройки функции <<ДЗ-фф 1ст.>>\hfill\vspace{-0.5\baselineskip}}\label{dz:dz1fftbl1}\\ 
\hline
\rowcolor{gray!30}
Параметр (Параметр на ИЧМ) & Условное обозначение на схеме & Значение/ Диапазон & Единица измерения & Шаг \\ 
\hline
\endfirsthead
\caption*{\hspace{3pt}\emph{Продолжение таблицы \ref{dz:dz1fftbl1}\hfill\vspace{-0.5\baselineskip}}} \\ % сделано по ГОСТ 2.105 п.6.8.7
\hline
\rowcolor{gray!30}
Параметр (Параметр на ИЧМ) & Условное обозначение на схеме & Значение/ Диапазон & Единица измерения & Шаг \\ 
%\hline
\endhead
%\hline
\endfoot
%\hline
\endlastfoot
%==+t1*PPPDIS1|VV_LVDIS
\centering Ввод функции в работу (Ввод\_функции) & \centering SGF1 & \centering 0 = Не предусмотрено\\1 = Предусмотрено & \centering -- & \centering \arraybackslash -- \\
\hline
\centering Угол наклона верхней стороны характеристики срабатывания для компенсации нагрузки (Фкн) & \centering -- & \centering -45,0 ... 0,0 & \centering град. & \centering \arraybackslash 0,1 \\
\hline
\centering Активное сопротивление срабатывания (Rср) & \centering -- & \centering 0,01 ... 500,00 & \centering Ом & \centering \arraybackslash 0,01 \\
\hline
\centering Реактивное сопротивление срабатывания (Xср) & \centering -- & \centering 0,01 ... 500,00 & \centering Ом & \centering \arraybackslash 0,01 \\
\hline
\centering Угол наклона характеристики срабатывания (Фхс) & \centering -- & \centering 30,0 ... 89,9 & \centering град. & \centering \arraybackslash 0,1 \\
\hline
\centering Режим направленности (Направленность) & \centering SGF3 & \centering 1 = Ненаправленная\\2 = Прямонаправленная\\3 = Обратнонаправленная & \centering -- & \centering \arraybackslash -- \\
\hline
\centering Подхват от ненаправленного пуска или отдельной ступени  (Подхват) & \centering SGF9 & \centering 0 = Не предусмотрено\\1 = Предусмотрено & \centering -- & \centering \arraybackslash -- \\
\hline
\centering Режим контроля от БК (Реж\_БК) & \centering SGF7 & \centering 1 = С контролем от БК-б\\2 = С контролем от БК-м & \centering -- & \centering \arraybackslash -- \\
\hline
\centering Режим контроля от БНН (Реж\_БНН) & \centering SGF8 & \centering 0 = Без контроля\\1 = С контролем & \centering -- & \centering \arraybackslash -- \\
\hline
\centering Режим контроля от ОПО (Реж\_ОПО) & \centering SGF6 & \centering 1 = По току\\2 = По току и напряжению\\3 = По току или по току и напряжению & \centering -- & \centering \arraybackslash -- \\
\hline
\centering Тип ПО (Контр\_ПО) & \centering SGF5 & \centering 1 = Без контроля\\2 = С контролем от БК\\3 = С контролем от ОПО & \centering -- & \centering \arraybackslash -- \\
\hline
\centering Вывод направленности при АУ (Ненапр\_АУ) & \centering SGF4 & \centering 0 = Не предусмотрено\\1 = Предусмотрено & \centering -- & \centering \arraybackslash -- \\
\hline
\centering Выдержка времени срабатывания (Tср) & \centering T1 & \centering 0,00 ... 10,00 & \centering с & \centering \arraybackslash 0,01 \\
\hline
%===t1
\end{longtable}
\normalsize

\item
Алгоритм работы второй ступени ДЗ по контуру <<фаза-фаза>> (<<ДЗ-фф 2ст.>>) выполнен в соответствии с рисунком \ref{dz:dz2ffimg1}. В таблице \ref{dz:dz2fftbl1} приведены параметры, необходимые для настройки функции <<ДЗ-фф 2ст.>>.

\small
\begin{longtable}{|>{\centering\arraybackslash}m{5.3cm}|>{\centering\arraybackslash}m{3.3cm}|>{\centering\arraybackslash}m{4.2cm}|>{\centering\arraybackslash}m{1.8cm}|>{\centering\arraybackslash}m{1cm}|}
\caption{Параметры для настройки функции <<ДЗ-фф 2ст.>>\hfill\vspace{-0.5\baselineskip}}\label{dz:dz2fftbl1}\\ 
\hline
\rowcolor{gray!30}
Параметр (Параметр на ИЧМ) & Условное обозначение на схеме & Значение/ Диапазон & Единица измерения & Шаг \\ 
\hline
\endfirsthead
%\caption{\emph{Продолжение\hfill\vspace{-0.5\baselineskip}}} \\ % Отступ обозначения таблицы от таблицы и выравнивание надписи слева
\caption*{\hspace{3pt}\emph{Продолжение таблицы \ref{dz:dz2fftbl1}\hfill\vspace{-0.5\baselineskip}}} \\ % сделано по ГОСТ 2.105 п.6.8.7
\hline
\rowcolor{gray!30}
Параметр (Параметр на ИЧМ) & Условное обозначение на схеме & Значение/ Диапазон & Единица измерения & Шаг \\ 
%\hline
\endhead
%\hline
\endfoot
%\hline
\endlastfoot
%==+t1*PPPDIS2|VV_LVDIS
\centering Ввод функции в работу (Ввод\_функции) & \centering SGF1 & \centering 0 = Не предусмотрено\\1 = Предусмотрено & \centering -- & \centering \arraybackslash -- \\
\hline
\centering Угол наклона верхней стороны характеристики срабатывания для компенсации нагрузки (Фкн) & \centering -- & \centering -45,0 ... 0,0 & \centering град. & \centering \arraybackslash 0,1 \\
\hline
\centering Активное сопротивление срабатывания (Rср) & \centering -- & \centering 0,01 ... 500,00 & \centering Ом & \centering \arraybackslash 0,01 \\
\hline
\centering Реактивное сопротивление срабатывания (Xср) & \centering -- & \centering 0,01 ... 500,00 & \centering Ом & \centering \arraybackslash 0,01 \\
\hline
\centering Угол наклона характеристики срабатывания (Фхс) & \centering -- & \centering 30,0 ... 89,9 & \centering град. & \centering \arraybackslash 0,1 \\
\hline
\centering Режим направленности (Направленность) & \centering SGF3 & \centering 1 = Ненаправленная\\2 = Прямонаправленная\\3 = Обратнонаправленная & \centering -- & \centering \arraybackslash -- \\
\hline
\centering Режим контроля от БК (Реж\_БК) & \centering SGF7 & \centering 1 = С контролем от БК-б\\2 = С контролем от БК-м & \centering -- & \centering \arraybackslash -- \\
\hline
\centering Режим контроля от БНН (Реж\_БНН) & \centering SGF8 & \centering 0 = Без контроля\\1 = С контролем & \centering -- & \centering \arraybackslash -- \\
\hline
\centering Режим контроля от ОПО (Реж\_ОПО) & \centering SGF6 & \centering 1 = По току\\2 = По току и напряжению\\3 = По току или по току и напряжению & \centering -- & \centering \arraybackslash -- \\
\hline
\centering Тип ПО (Контр\_ПО) & \centering SGF5 & \centering 1 = Без контроля\\2 = С контролем от БК\\3 = С контролем от ОПО & \centering -- & \centering \arraybackslash -- \\
\hline
\centering Вывод направленности при АУ (Ненапр\_АУ) & \centering SGF4 & \centering 0 = Не предусмотрено\\1 = Предусмотрено & \centering -- & \centering \arraybackslash -- \\
\hline
\centering Выдержка времени срабатывания (Tср) & \centering T1 & \centering 0,00 ... 10,00 & \centering с & \centering \arraybackslash 0,01 \\
\hline
%===t1
\end{longtable}
\normalsize

\vspace{3mm}
\begin{figure}[H]
\centering
\includegraphics[width=0.9\textwidth,height=0.9\textheight,keepaspectratio]{\fbpath/301. ДЗ ВВ/_latex/img/dz2ff.pdf}
\captionof{figure}{Функциональная схема <<ДЗ-фф 2ст.>>}\label{dz:dz2ffimg1}
\end{figure}

\item
Алгоритм работы третьей ступени ДЗ по контуру <<фаза-фаза>> (<<ДЗ-фф 3ст.>>) выполнен в соответствии с рисунком \ref{dz:dz3ffimg1}. В таблице \ref{dz:dz3fftbl1} приведены параметры, необходимые для настройки функции <<ДЗ-фф 3ст.>>.

\small
\begin{longtable}{|>{\centering\arraybackslash}m{5.3cm}|>{\centering\arraybackslash}m{3.3cm}|>{\centering\arraybackslash}m{4.2cm}|>{\centering\arraybackslash}m{1.8cm}|>{\centering\arraybackslash}m{1cm}|}
\caption{Параметры для настройки функции <<ДЗ-фф 3ст.>>\hfill\vspace{-0.5\baselineskip}}\label{dz:dz3fftbl1}\\ 
\hline
\rowcolor{gray!30}
Параметр (Параметр на ИЧМ) & Условное обозначение на схеме & Значение/ Диапазон & Единица измерения & Шаг \\ 
\hline
\endfirsthead
\caption*{\hspace{3pt}\emph{Продолжение таблицы \ref{dz:dz3fftbl1}\hfill\vspace{-0.5\baselineskip}}} \\ % сделано по ГОСТ 2.105 п.6.8.7
\hline
\rowcolor{gray!30}
Параметр (Параметр на ИЧМ) & Условное обозначение на схеме & Значение/ Диапазон & Единица измерения & Шаг \\ 
%\hline
\endhead
%\hline
\endfoot
%\hline
\endlastfoot
%==+t1*PPPDIS3|VV_LVDIS
\centering Ввод функции в работу (Ввод\_функции) & \centering SGF1 & \centering 0 = Не предусмотрено\\1 = Предусмотрено & \centering -- & \centering \arraybackslash -- \\
\hline
\centering Угол наклона верхней стороны характеристики срабатывания для компенсации нагрузки (Фкн) & \centering -- & \centering -45,0 ... 0,0 & \centering град. & \centering \arraybackslash 0,1 \\
\hline
\centering Активное сопротивление срабатывания (Rср) & \centering -- & \centering 0,01 ... 500,00 & \centering Ом & \centering \arraybackslash 0,01 \\
\hline
\centering Реактивное сопротивление срабатывания (Xср) & \centering -- & \centering 0,01 ... 500,00 & \centering Ом & \centering \arraybackslash 0,01 \\
\hline
\centering Угол наклона характеристики срабатывания (Фхс) & \centering -- & \centering 30,0 ... 89,9 & \centering град. & \centering \arraybackslash 0,1 \\
\hline
\centering Режим направленности (Направленность) & \centering SGF3 & \centering 1 = Ненаправленная\\2 = Прямонаправленная\\3 = Обратнонаправленная & \centering -- & \centering \arraybackslash -- \\
\hline
\centering Режим контроля от БК (Реж\_БК) & \centering SGF7 & \centering 1 = С контролем от БК-б\\2 = С контролем от БК-м & \centering -- & \centering \arraybackslash -- \\
\hline
\centering Режим контроля от БНН (Реж\_БНН) & \centering SGF8 & \centering 0 = Без контроля\\1 = С контролем & \centering -- & \centering \arraybackslash -- \\
\hline
\centering Режим контроля от ОПО (Реж\_ОПО) & \centering SGF6 & \centering 1 = По току\\2 = По току и напряжению\\3 = По току или по току и напряжению & \centering -- & \centering \arraybackslash -- \\
\hline
\centering Тип ПО (Контр\_ПО) & \centering SGF5 & \centering 1 = Без контроля\\2 = С контролем от БК\\3 = С контролем от ОПО & \centering -- & \centering \arraybackslash -- \\
\hline
\centering Вывод направленности при АУ (Ненапр\_АУ) & \centering SGF4 & \centering 0 = Не предусмотрено\\1 = Предусмотрено & \centering -- & \centering \arraybackslash -- \\
\hline
\centering Выдержка времени срабатывания (Tср) & \centering T1 & \centering 0,00 ... 10,00 & \centering с & \centering \arraybackslash 0,01 \\
\hline
%===t1
\end{longtable}
\normalsize

\vspace{3mm}
\begin{figure}[H]
\centering
\includegraphics[width=0.9\textwidth,height=0.9\textheight,keepaspectratio]{\fbpath/301. ДЗ ВВ/_latex/img/dz3ff.pdf}
\captionof{figure}{Функциональная схема <<ДЗ-фф 3ст.>>}\label{dz:dz3ffimg1}
\end{figure}

\end{enumerate}

%%%%%%%%%%%%%%%%%%%%%%%%%%%%%%%%%%%%%%%%%%%%%%%%%%%%%%%%%%%%%%%%%%%%%% ОРГАН НАПРАВЛЕНИЯ %%%%%%%%%%%%%%%%%%%%%%%%%%%%%%%%%%%%%%%%%%%%%%%%%%
\item\label{sec:on} Орган направления (ОН)
\begin{enumerate}[label=\arabic{section}.\arabic{subsection}.\arabic{enumi}.\arabic*, labelsep=4pt, leftmargin=0em, itemindent=65pt, parsep=0pt]

\item
В кольцевых сетях, в сетях с двухсторонним питанием и особенно в сложных сетях с несколькими источниками питания в большинстве случаев ненаправленная защита не может обеспечить селективность 
действия. Для  селективной  работы в таких случаях необходимо применение органа  направления мощности совместно с измерительным токовым органом защиты. Защита действует на отключение только 
в  том  случае, если срабатывает как токовый измерительный орган, так и орган  направления, который разрешает действие при направлении мощности КЗ от шин в линию.
\item
Функция  <<ОН>>  становится  активна,  если  введена  в  работу  одна  из  ступеней  ДЗ и в ней настроен контроль направленности. 
\item
ОН  применим  для  реализации  направленных  ступеней  ДЗ.  Направленная  ступень  ДЗ  будет действовать  только  в  том  случае,  когда  сработает  реле  сопротивления,  а  также  ОН,  разрешающий работать ДЗ в прямом направлении. Функция ОН содержит в себе три органа, определяющих направление. 
\item
ОН  выполнен  в  виде  отдельного  функционального  модуля  с  собственными  уставками.  ОН  ДЗ осуществляет построение характеристики ОН по уставкам, ограничивающим характеристику срабатывания органа в прямом направлении во втором и четвертом квадрантах (<<Угол во 2 квадр.>> - $\varphi2$ и <<Угол в 4 квадр.>> – $\varphi1$ соответственно), и, используя положение  вектора  сопротивления  на  плоскости  характеристики, формирует выходные сигналы прямого либо обратного направления согласно характеристике, изображенной на \ref{dz:onimg0}.

\vspace{3mm}
\begin{figure}[H]
\centering
\includegraphics[width=0.5\textwidth,height=0.5\textheight,keepaspectratio]{\fbpath/301. ДЗ ВВ/_latex/img/charon.pdf}
\captionof{figure}{{Характеристика ОН}}\label{dz:onimg0}
\end{figure}

\item
ОН ДЗ может работать в одном из режимов направленности:
\begin{itemize} 
    \item в прямом направлении, если вектор сопротивления находится в диапазоне от $\varphi1$ до $\varphi2$; 
    \item в обратном направлении, если вектор сопротивления находится в диапазоне от $\varphi1-180^\circ$ до $\varphi2-180^\circ$; 
\item без направления, если вектор сопротивления находится в области нечувствительности.
\end{itemize}

\item
Строгая направленность обеспечивается действиями ступеней ДЗ при КЗ вблизи установки защиты используется специальный ОН на аварийных составляющих, который вводится при снижении линейного напряжения ниже 5 В.
\item
Данный орган контролирует направление повреждения и блокирует срабатывание ступеней ДЗ при КЗ <<за спиной>>. ОН на аварийных составляющих выполнен таким образом, что обеспечивает правильное 
определение  направления  повреждения  при  любых  видах  КЗ  (в  том  числе  при  близких  с  просадкой напряжения).  Обеспечивается  отсутствие  ложных  срабатываний  ОН  при  КЗ  <<за  спиной>>  при  токах  до $20 \cdot I_{\textsl{ном}}$.
\item
Отказ ОН возможен при включении на КЗ, если защита подключена к ТН, установленного на линии. Эта  особенность  характерна  всем  известным  способам,  работа  которых  опирается  на  напряжение 
предшествующего режима.

\item
Принцип работы линейных органов следующий:
\begin{itemize} 
\item в  качестве  воздействующих  величин  используются  действующие  значения  первой гармоники линейных токов и значения линейных напряжений поляризации рассчитанные по формулам (\ref{eq:group1}); 
\item разрешающими  условиями  для  работы  направленного  органа  являются  превышение  тока значения $0,04 \cdot I_{\textsl{ном}}$ и превышение напряжения значения 1 вольт; 
\item фаза вектора сопротивления, рассчитанного по формулам (\ref{eq:group2}) находится в области срабатывания (прямое или обратное направление).
\end{itemize}

\item
В случае снижения напряжения ниже 0,9 вольт орган направления в течение 100 мс продолжит работу по памяти, и по истечении данного времени произойдет подхват сигнала направления. 

\item
Алгоритм работы органа направления (<<ОН>>) выполнен в соответствии с рисунком \ref{dz:onimg1}. Алгоритм работы логики определения направления, входящий в состав органа направления, выполнен в соответствии с рисунком \ref{dz:onimg2}. В таблице \ref{dz:ontbl1} приведены параметры, необходимые для настройки функции <<ОН>>.

\vspace{3mm}
\begin{figure}[H]
\centering
\includegraphics[width=1\textwidth,height=1\textheight,keepaspectratio]{\fbpath/301. ДЗ ВВ/_latex/img/on1.pdf}
\captionof{figure}{Функциональная схема логики определения направления}\label{dz:onimg2}
\end{figure}

\small
\begin{longtable}{|>{\centering\arraybackslash}m{5.3cm}|>{\centering\arraybackslash}m{3.3cm}|>{\centering\arraybackslash}m{4.2cm}|>{\centering\arraybackslash}m{1.8cm}|>{\centering\arraybackslash}m{1cm}|}
\caption{Параметры для настройки функции <<ОН>>\hfill\vspace{-0.5\baselineskip}}\label{dz:ontbl1}\\ 
\hline
\rowcolor{gray!30}
Параметр (Параметр на ИЧМ) & Условное обозначение на схеме & Значение/ Диапазон & Единица измерения & Шаг \\ 
\hline
\endfirsthead
\caption*{\hspace{3pt}\emph{Продолжение таблицы \ref{dz:ontbl1}\hfill\vspace{-0.5\baselineskip}}} \\ % сделано по ГОСТ 2.105 п.6.8.7
\hline
\rowcolor{gray!30}
Параметр (Параметр на ИЧМ) & Условное обозначение на схеме & Значение/ Диапазон & Единица измерения & Шаг \\ 
%\hline
\endhead
%\hline
\endfoot
%\hline
\endlastfoot
%==+t1*DISRDIR1|VV_LVDIS
\centering Угол направленности во II квадранте (Фнапр\_II) & \centering φ2 & \centering 90,0 ... 130,0 & \centering град. & \centering \arraybackslash 0,1 \\
\hline
\centering Угол направленности в IV квадранте (Фнапр\_IV) & \centering φ1 & \centering -40,0 ... 0,0 & \centering град. & \centering \arraybackslash 0,1 \\
\hline
%===t1
\end{longtable}
\normalsize

\vspace{3mm}
\begin{figure}[H]
\centering
\includegraphics[width=0.8\textwidth,height=0.8\textheight,keepaspectratio]{\fbpath/301. ДЗ ВВ/_latex/img/on2.pdf}
\captionof{figure}{Функциональная схема <<ОН>>}\label{dz:onimg1}
\end{figure}

\end{enumerate}

%%%%%%%%%%%%%%%%%%%%%%%%%%%%%%%%%%%%%%%%%%%%%%%%%%%%%%%%%%%%%%%%%%%%%% Общие пусковые органы ДЗ (ОПО ДЗ) %%%%%%%%%%%%%%%%%%%%%%%%%%%%%%%%%%%%%%%%%%%%%%%%%%
\item\label{sec:opo} Общие пусковые органы ДЗ (ОПО ДЗ)
\begin{enumerate}[label=\arabic{section}.\arabic{subsection}.\arabic{enumi}.\arabic*, labelsep=4pt, leftmargin=0em, itemindent=65pt, parsep=0pt]

\item
Чтобы  избежать  ложного  срабатывания  ступени  дистанционной  защиты  (ДЗ),  осуществляется дополнительный  контроль  срабатывания  пусковых  органов  дистанционной  защиты  несколькими 
способами:
\begin{itemize}
    \item с пуском по току; 
    \item с пуском по току и напряжению; 
    \item с пуском по току или по току и напряжению.
\end{itemize}
\item
Функция становится активна, если введена в работу одна из ступеней ДЗ и в ней выбран контроль от общих пусковых органов (режим  SGF5 -- <<С контролем от ОПО>> ).
\item
Сигнал общего пуска, пуска от измерительного органа ОПО  ДЗ по току формируется, если максимальное действующее значение первой гармоники одного из фазных токов превысило уставку.
\item
Сигнал общего пуска, пуска от измерительного органа ОПО ДЗ по току и напряжению формируется при снижении одного из линейных напряжений ниже значения уставки и превышении одним из фазных токов значения выше уставки.

\item
Алгоритм работы функции ОПО ДЗ (<<ОПО ДЗ>>) выполнен в соответствии с рисунком \ref{dz:opoimg1}. В таблице \ref{dz:opotbl1} приведены параметры, необходимые для настройки функции <<ОПО ДЗ>>.

\vspace{3mm}
\begin{figure}[H]
\centering
\includegraphics[width=0.8\textwidth,height=0.8\textheight,keepaspectratio]{\fbpath/301. ДЗ ВВ/_latex/img/opo.pdf}
\captionof{figure}{Функциональная схема <<ОПО ДЗ>>}\label{dz:opoimg1}
\end{figure}

\small
\begin{longtable}{|>{\centering\arraybackslash}m{5.3cm}|>{\centering\arraybackslash}m{3.3cm}|>{\centering\arraybackslash}m{4.2cm}|>{\centering\arraybackslash}m{1.8cm}|>{\centering\arraybackslash}m{1cm}|}
\caption{Параметры для настройки функции <<ОПО ДЗ>>\hfill\vspace{-0.5\baselineskip}}\label{dz:opotbl1}\\ 
\hline
\rowcolor{gray!30}
Параметр (Параметр на ИЧМ) & Условное обозначение на схеме & Значение/ Диапазон & Единица измерения & Шаг \\ 
\hline
\endfirsthead
\caption*{\hspace{3pt}\emph{Продолжение таблицы \ref{dz:opotbl1}\hfill\vspace{-0.5\baselineskip}}} \\ % сделано по ГОСТ 2.105 п.6.8.7
\hline
\rowcolor{gray!30}
Параметр (Параметр на ИЧМ) & Условное обозначение на схеме & Значение/ Диапазон & Единица измерения & Шаг \\ 
%\hline
\endhead
%\hline
\endfoot
%\hline
\endlastfoot
%==+t1*DISPGSU1|VV_LVDIS
\centering Ток срабатывания ПО ДЗ по току (Iср) & \centering I\verb|>>| & \centering 0,10 ... 5,00 & \centering о.е. & \centering \arraybackslash 0,01 \\
\hline
\centering Ток срабатывания ПО ДЗ по току и напряжению (IсрU) & \centering I> & \centering 0,10 ... 5,00 & \centering о.е. & \centering \arraybackslash 0,01 \\
\hline
\centering Напряжение срабатывания ПО по току и напряжению (Uср) & \centering U< & \centering 2,0 ... 100,0 & \centering В & \centering \arraybackslash 0,1 \\
\hline
\centering Режим контроля от БНН (Реж\_БНН) & \centering SGF1 & \centering 0 = Не предусмотрено\\1 = Предусмотрено & \centering -- & \centering \arraybackslash -- \\
\hline
%===t1
\end{longtable}
\normalsize
\end{enumerate}
%%%%%%%%%%%%%%%%%%%%%%%%%%%%%%%%%%%%%%%%%%%%%%%%%%%%%%%%%%%%%%%%%%%%%% КОНЕЦ Общие пусковые органы ДЗ (ОПО ДЗ) %%%%%%%%%%%%%%%%%%%%%%%%%%%%%%%%%%%%%%%%%%%%%%%%%%

%%%%%%%%%%%%%%%%%%%%%%%%%%%%%%%%%%%%%%%%%%%%%%%%%%%%%%%%%%%%%%%%%%%%%%  Автоматическое ускорение ДЗ (АУ ДЗ)  %%%%%%%%%%%%%%%%%%%%%%%%%%%%%%%%%%%%%%%%%%%%%%%%%%
\item\label{sec:audz} Автоматическое ускорение ДЗ (АУ ДЗ)
\begin{enumerate}[label=\arabic{section}.\arabic{subsection}.\arabic{enumi}.\arabic*, labelsep=4pt, leftmargin=0em, itemindent=65pt, parsep=0pt]

\item    
В  устройстве предусмотрен режим автоматического ускорения ступеней ДЗ.  Выбор ускоряемых ступеней защиты осуществляется из соображений надежного охвата, защищаемого энергообъекта.
\item 
С  помощью  уставки  SGF1 <<Выбор\_ст>>  функции АУ обеспечивается  возможность  автоматического  ускоренного пуска от второй или третьей ступеней ДЗ. Пусковые условия автоматического ускорения формируются без выдержки времени.

\item
Алгоритм работы функции АУ ДЗ (<<АУ ДЗ>>) выполнен в соответствии с рисунком \ref{dz:audzimg1}. В таблице \ref{dz:audztbl1} приведены параметры, необходимые для настройки функции <<АУ ДЗ>>.

\vspace{3mm}
\begin{figure}[H]
\centering
\includegraphics[width=0.8\textwidth,height=0.8\textheight,keepaspectratio]{\fbpath/301. ДЗ ВВ/_latex/img/audz.pdf}
\captionof{figure}{Функциональная схема <<АУ ДЗ>>}\label{dz:audzimg1}
\end{figure}

\small
\begin{longtable}{|>{\centering\arraybackslash}m{5.3cm}|>{\centering\arraybackslash}m{3.3cm}|>{\centering\arraybackslash}m{4.2cm}|>{\centering\arraybackslash}m{1.8cm}|>{\centering\arraybackslash}m{1cm}|}
\caption{Параметры для настройки функции <<АУ ДЗ>>\hfill\vspace{-0.5\baselineskip}}\label{dz:audztbl1}\\ 
\hline
\rowcolor{gray!30}
Параметр (Параметр на ИЧМ) & Условное обозначение на схеме & Значение/ Диапазон & Единица измерения & Шаг \\ 
\hline
\endfirsthead
\caption*{\hspace{3pt}\emph{Продолжение таблицы \ref{dz:audztbl1}\hfill\vspace{-0.5\baselineskip}}} \\ % сделано по ГОСТ 2.105 п.6.8.7
\hline
\rowcolor{gray!30}
Параметр (Параметр на ИЧМ) & Условное обозначение на схеме & Значение/ Диапазон & Единица измерения & Шаг \\ 
%\hline
\endhead
%\hline
\endfoot
%\hline
\endlastfoot
%==+t1*AUAPSSU1|VV_LVDIS
\centering Выбор ускоряемой ступени (Выбор\_ст) & \centering SGF1 & \centering 1 = 2 ступень\\2 = 3 ступень & \centering -- & \centering \arraybackslash -- \\
\hline
%===t1
\end{longtable}
\normalsize

\end{enumerate}
%%%%%%%%%%%%%%%%%%%%%%%%%%%%%%%%%%%%%%%%%%%%%%%%%%%%%%%%%%%%%%%%%% КОНЕЦ  Автоматическое ускорение ДЗ (АУ ДЗ) %%%%%%%%%%%%%%%%%%%%%%%%%%%%%%%%%%%%%%%%%%%%%%%%%%

%%%%%%%%%%%%%%%%%%%%%%%%%%%%%%%%%%%%%%%%%%%%%%%%%%%%%%%%%%%%%%%%%%%%%%  Оперативное ускорение ДЗ (ОУ ДЗ)  %%%%%%%%%%%%%%%%%%%%%%%%%%%%%%%%%%%%%%%%%%%%%%%%%%
\item\label{sec:oudz} Оперативное ускорение ДЗ (ОУ ДЗ)
\begin{enumerate}[label=\arabic{section}.\arabic{subsection}.\arabic{enumi}.\arabic*, labelsep=4pt, leftmargin=0em, itemindent=65pt, parsep=0pt]

\item   
Оперативное ускорение ДЗ может вводиться при:
\begin{itemize}
    \item выводе основной защиты оборудования; 
    \item опробовании оборудования рабочим напряжением; 
    \item операциях с разъединителями на противоположном конце ВЛ. 
\end{itemize}

\item
В  данном  алгоритме  при  введенном  в  работу  ключе  <<Ввод  ОУ  ДЗ>>  и  положением  уставки  SGF1 предусмотрена возможность выбора от пуска какой ступени ДЗ будет формироваться срабатывание ОУ ДЗ. 

\item
Срабатывание оперативного ускорения от пуска ступеней ДЗ происходит по истечению выдержки времени таймера Т1. 

\item
Алгоритм работы функции ОУ ДЗ (<<ОУ ДЗ>>) выполнен в соответствии с рисунком \ref{dz:oudzimg1}. В таблице \ref{dz:oudztbl1} приведены параметры, необходимые для настройки функции <<ОУ ДЗ>>.

\vspace{3mm}
\begin{figure}[H]
\centering
\includegraphics[width=0.8\textwidth,height=0.8\textheight,keepaspectratio]{\fbpath/301. ДЗ ВВ/_latex/img/oudz.pdf}
\captionof{figure}{Функциональная схема <<ОУ ДЗ>>}\label{dz:oudzimg1}
\end{figure}

\small
\begin{longtable}{|>{\centering\arraybackslash}m{5.3cm}|>{\centering\arraybackslash}m{3.3cm}|>{\centering\arraybackslash}m{4.2cm}|>{\centering\arraybackslash}m{1.8cm}|>{\centering\arraybackslash}m{1cm}|}
\caption{Параметры для настройки функции <<ОУ ДЗ>>\hfill\vspace{-0.5\baselineskip}}\label{dz:oudztbl1}\\ 
\hline
\rowcolor{gray!30}
Параметр (Параметр на ИЧМ) & Условное обозначение на схеме & Значение/ Диапазон & Единица измерения & Шаг \\ 
\hline
\endfirsthead
\caption*{\hspace{3pt}\emph{Продолжение таблицы \ref{dz:oudztbl1}\hfill\vspace{-0.5\baselineskip}}} \\ % сделано по ГОСТ 2.105 п.6.8.7
\hline
\rowcolor{gray!30}
Параметр (Параметр на ИЧМ) & Условное обозначение на схеме & Значение/ Диапазон & Единица измерения & Шаг \\ 
%\hline
\endhead
%\hline
\endfoot
%\hline
\endlastfoot
%==+t1*OPAPSSU1|VV_LVDIS
\centering Выбор оперативно ускоряемой ступени (Выбор\_ст\_ОУ) & \centering SGF1 & \centering 1 = 2 ступень\\2 = 3 ступень & \centering -- & \centering \arraybackslash -- \\
\hline
\centering Выдержка времени срабатывания оперативного ускорения (T\_ОУ) & \centering T1 & \centering 0,00 ... 3,00 & \centering с & \centering \arraybackslash 0,01 \\
\hline
%===t1
\end{longtable}
\normalsize

\end{enumerate}
%%%%%%%%%%%%%%%%%%%%%%%%%%%%%%%%%%%%%%%%%%%%%%%%%%%%%%%%%%%%%%%%%%%%%% КОНЕЦ  Оперативное ускорение ДЗ (ОУ ДЗ) %%%%%%%%%%%%%%%%%%%%%%%%%%%%%%%%%%%%%%%%%%%%%%%%%%

%%%%%%%%%%%%%%%%%%%%%%%%%%%%%%%%%%%%%%%%%%%%%%%%%%%%%%%%%%%%%%%%%%%%%%  Внешняя логика функции ДЗ   %%%%%%%%%%%%%%%%%%%%%%%%%%%%%%%%%%%%%%%%%%%%%%%%%%
\item\label{sec:vldz} Внешняя логика функции ДЗ
\begin{enumerate}[label=\arabic{section}.\arabic{subsection}.\arabic{enumi}.\arabic*, labelsep=4pt, leftmargin=0em, itemindent=65pt, parsep=0pt]

\item
Для функции предусмотрен общий выходной сигнал «ДЗ: Пуск», который формируется при наличии пусковых сигналов от одной из ступеней ДЗ. 

\item
Алгоритм работы  внешней логики ДЗ выполнен в соответствии с рисунком \ref{dz:vlimg1}.

\vspace{3mm}
\begin{figure}[H]
\centering
\includegraphics[width=0.6\textwidth,height=0.6\textheight,keepaspectratio]{\fbpath/301. ДЗ ВВ/_latex/img/vldz.pdf}
\captionof{figure}{Функциональная схема внешней логики ДЗ}\label{dz:vlimg1}
\end{figure}

\end{enumerate}
%%%%%%%%%%%%%%%%%%%%%%%%%%%%%%%%%%%%%%%%%%%%%%%%%%%%%%%%%%%%%%%%%%%%%% КОНЕЦ  Внешняя логика функции ДЗ %%%%%%%%%%%%%%%%%%%%%%%%%%%%%%%%%%%%%%%%%%%%%%%%%%


\end{enumerate}
